% --- Archon physics formalization source begin ---
% archon:physics
% archon:covers PhyXMiniProblems/problem_phyx_mini_0385.lean
% archon:source-report reports/phyx_mini/problem_phyx_mini_0385.source.json

\chapter{Physics problem phyx\_mini\_0385}
\label{ch:PhyXMiniProblems_problem_phyx_mini_0385}

\paragraph{Problem source.}
\#\# Physical scenario

A 2.5-m-tall steel cylinder has a cross-sectional area of 1.5 \textbackslash{}m\textasciicircum{}2\textbackslash{}. At the bottom, with a height of 0.5 m, is liquid water, on top of which is a 1-m-high layer of gasoline. This is shown in the figure. The gasoline surface is exposed to atmospheric air at 101 kPa.

\#\# Question

What is the highest pressure in the water?

\#\# Answer choices

- A: A: 490kPa

- B: B: 113.2kPa

- C: C: 154kPa

- D: D: 10.2kPa

\#\# Figure caption (auxiliary; use the image as primary evidence)

The image depicts a cross-sectional view of a vertical container. It contains three distinct layers:

1. **Bottom Layer**: 
   - Labeled "H2O" (water).
   - Occupies a height of 0.5 meters from the bottom.

2. **Middle Layer**: 
   - Labeled "Gasoline".
   - Sits above the water layer.
   - Has a height of 1 meter.

3. **Top Layer**: 
   - Labeled "Air".
   - Fills the remaining space above the gasoline.
   - The entire height from the bottom of the container to the top is 2.5 meters, indicating the air layer occupies the remaining space.

The pressure at the top is indicated as \textbackslash{}( P\_0 \textbackslash{}).

The container's sides are straight and vertical, and each layer is clearly separated and labeled in the image.

\#\# Dataset metadata

Ideal Gases and Kinetic Theory; Physical Model Grounding Reasoning, Multi-Formula Reasoning

\paragraph{Recorded answer/context.}
B: B: 113.2kPa

\paragraph{Figure/image path.}
/root/proposal\_for\_physic/science-mango/phyx\_mini\_run/phyx\_data/test\_image/385.png

\paragraph{Formalization target.}
create a compiling Lean file with sorry bodies at `PhyXMiniProblems/problem\_phyx\_mini\_0385.lean`. The Lean declarations must preserve the physical quantities, dimensions or dimensional roles, figure labels, governing-law hypotheses, and final relation expressed by this problem.
Use Mathlib/Physlib names found through LeanExplore where available. If a physics API is missing, introduce faithful local abstractions rather than scalar placeholder aliases.

\begin{theorem}[Physics formalization target]
\label{thm:physics:phyx_mini_0385:target}
\lean{PhyXMiniProblems.ProblemPhyXMini0385.highestPressureInWater_fromSource}
\uses{def:physics:phyx-mini-0385:phyxminiproblems-problemphyxmini0385-lengthinmeters, def:physics:phyx-mini-0385:phyxminiproblems-problemphyxmini0385-pressureinpascals, def:physics:phyx-mini-0385:phyxminiproblems-problemphyxmini0385-pressureinkilopascals, def:physics:phyx-mini-0385:phyxminiproblems-problemphyxmini0385-densityinkilogramspercubicmeter, def:physics:phyx-mini-0385:phyxminiproblems-problemphyxmini0385-accelerationinmeterspersecondsquared, def:physics:phyx-mini-0385:phyxminiproblems-problemphyxmini0385-layeredcylindersetup, def:physics:phyx-mini-0385:phyxminiproblems-problemphyxmini0385-matcheslayeredcylinderscenario, def:physics:phyx-mini-0385:phyxminiproblems-problemphyxmini0385-matchesproblemreadouts, def:physics:phyx-mini-0385:phyxminiproblems-problemphyxmini0385-matchessuppliedfigure, def:physics:phyx-mini-0385:phyxminiproblems-problemphyxmini0385-usesstandarddensityandgravitycalibration, def:physics:phyx-mini-0385:phyxminiproblems-problemphyxmini0385-hasphysicallayerparameters, def:physics:phyx-mini-0385:phyxminiproblems-problemphyxmini0385-satisfieslayeredhydrostaticlaws, def:physics:phyx-mini-0385:phyxminiproblems-problemphyxmini0385-ishighestpressureinwater, def:physics:phyx-mini-0385:phyxminiproblems-problemphyxmini0385-isuniqueclosestdisplayedpressure}
The assigned autoformalize agent should translate this physics problem into Lean declarations in the covered file, with theorem and lemma proof bodies written as `by sorry`.
\end{theorem}
\begin{proof}
This is an autoformalization task, not a proof task. Produce statements that can later be proved by the physics prover without weakening the physical model.
\end{proof}

% --- Archon declaration topology begin ---
\section{Declaration topology}
\begin{definition}[mass Density Dimension]
  \label{def:physics:phyx-mini-0385:phyxminiproblems-problemphyxmini0385-massdensitydimension}
  \lean{PhyXMiniProblems.ProblemPhyXMini0385.massDensityDimension}
  The physical dimension of mass density, M L⁻³
\end{definition}
\begin{proof}
  This is the declared model definition of mass Density Dimension.
\end{proof}
\begin{definition}[acceleration Dimension]
  \label{def:physics:phyx-mini-0385:phyxminiproblems-problemphyxmini0385-accelerationdimension}
  \lean{PhyXMiniProblems.ProblemPhyXMini0385.accelerationDimension}
  The physical dimension of acceleration, L T⁻²
\end{definition}
\begin{proof}
  This is the declared model definition of acceleration Dimension.
\end{proof}
\begin{definition}[Length Quantity]
  \label{def:physics:phyx-mini-0385:phyxminiproblems-problemphyxmini0385-lengthquantity}
  \lean{PhyXMiniProblems.ProblemPhyXMini0385.LengthQuantity}
  A nonnegative physical length independent of the readout unit
\end{definition}
\begin{proof}
  This is the declared model definition of Length Quantity.
\end{proof}
\begin{definition}[Area Quantity]
  \label{def:physics:phyx-mini-0385:phyxminiproblems-problemphyxmini0385-areaquantity}
  \lean{PhyXMiniProblems.ProblemPhyXMini0385.AreaQuantity}
  A nonnegative physical cross-sectional area
\end{definition}
\begin{proof}
  This is the declared model definition of Area Quantity.
\end{proof}
\begin{definition}[Mass Density Quantity]
  \label{def:physics:phyx-mini-0385:phyxminiproblems-problemphyxmini0385-massdensityquantity}
  \lean{PhyXMiniProblems.ProblemPhyXMini0385.MassDensityQuantity}
  \uses{def:physics:phyx-mini-0385:phyxminiproblems-problemphyxmini0385-massdensitydimension}
  A nonnegative physical mass density
\end{definition}
\begin{proof}
  This is the declared model definition of Mass Density Quantity.
\end{proof}
\begin{definition}[Acceleration Quantity]
  \label{def:physics:phyx-mini-0385:phyxminiproblems-problemphyxmini0385-accelerationquantity}
  \lean{PhyXMiniProblems.ProblemPhyXMini0385.AccelerationQuantity}
  \uses{def:physics:phyx-mini-0385:phyxminiproblems-problemphyxmini0385-accelerationdimension}
  A nonnegative physical acceleration magnitude
\end{definition}
\begin{proof}
  This is the declared model definition of Acceleration Quantity.
\end{proof}
\begin{definition}[length Readout]
  \label{def:physics:phyx-mini-0385:phyxminiproblems-problemphyxmini0385-lengthreadout}
  \lean{PhyXMiniProblems.ProblemPhyXMini0385.lengthReadout}
  \uses{def:physics:phyx-mini-0385:phyxminiproblems-problemphyxmini0385-lengthquantity}
  Read a physical length in a selected length unit
\end{definition}
\begin{proof}
  This is the declared model definition of length Readout.
\end{proof}
\begin{definition}[area Readout]
  \label{def:physics:phyx-mini-0385:phyxminiproblems-problemphyxmini0385-areareadout}
  \lean{PhyXMiniProblems.ProblemPhyXMini0385.areaReadout}
  \uses{def:physics:phyx-mini-0385:phyxminiproblems-problemphyxmini0385-areaquantity}
  Read a physical area in the square of a selected length unit
\end{definition}
\begin{proof}
  This is the declared model definition of area Readout.
\end{proof}
\begin{definition}[pressure Readout]
  \label{def:physics:phyx-mini-0385:phyxminiproblems-problemphyxmini0385-pressurereadout}
  \lean{PhyXMiniProblems.ProblemPhyXMini0385.pressureReadout}
  Read pressure in the coherent unit induced by selected base units
\end{definition}
\begin{proof}
  This is the declared model definition of pressure Readout.
\end{proof}
\begin{definition}[density Readout]
  \label{def:physics:phyx-mini-0385:phyxminiproblems-problemphyxmini0385-densityreadout}
  \lean{PhyXMiniProblems.ProblemPhyXMini0385.densityReadout}
  \uses{def:physics:phyx-mini-0385:phyxminiproblems-problemphyxmini0385-massdensityquantity}
  Read density in a selected mass unit per selected length unit cubed
\end{definition}
\begin{proof}
  This is the declared model definition of density Readout.
\end{proof}
\begin{definition}[acceleration Readout]
  \label{def:physics:phyx-mini-0385:phyxminiproblems-problemphyxmini0385-accelerationreadout}
  \lean{PhyXMiniProblems.ProblemPhyXMini0385.accelerationReadout}
  \uses{def:physics:phyx-mini-0385:phyxminiproblems-problemphyxmini0385-accelerationquantity}
  Read acceleration in selected length and time units
\end{definition}
\begin{proof}
  This is the declared model definition of acceleration Readout.
\end{proof}
\begin{definition}[length In Meters]
  \label{def:physics:phyx-mini-0385:phyxminiproblems-problemphyxmini0385-lengthinmeters}
  \lean{PhyXMiniProblems.ProblemPhyXMini0385.lengthInMeters}
  \uses{def:physics:phyx-mini-0385:phyxminiproblems-problemphyxmini0385-lengthquantity, def:physics:phyx-mini-0385:phyxminiproblems-problemphyxmini0385-lengthreadout}
  Meter readout of a physical length
\end{definition}
\begin{proof}
  This is the declared model definition of length In Meters.
\end{proof}
\begin{definition}[area In Square Meters]
  \label{def:physics:phyx-mini-0385:phyxminiproblems-problemphyxmini0385-areainsquaremeters}
  \lean{PhyXMiniProblems.ProblemPhyXMini0385.areaInSquareMeters}
  \uses{def:physics:phyx-mini-0385:phyxminiproblems-problemphyxmini0385-areaquantity, def:physics:phyx-mini-0385:phyxminiproblems-problemphyxmini0385-areareadout}
  Square-meter readout of a physical area
\end{definition}
\begin{proof}
  This is the declared model definition of area In Square Meters.
\end{proof}
\begin{definition}[pressure In Pascals]
  \label{def:physics:phyx-mini-0385:phyxminiproblems-problemphyxmini0385-pressureinpascals}
  \lean{PhyXMiniProblems.ProblemPhyXMini0385.pressureInPascals}
  \uses{def:physics:phyx-mini-0385:phyxminiproblems-problemphyxmini0385-pressurereadout}
  Pascal readout of an absolute pressure
\end{definition}
\begin{proof}
  This is the declared model definition of pressure In Pascals.
\end{proof}
\begin{definition}[pressure In Kilopascals]
  \label{def:physics:phyx-mini-0385:phyxminiproblems-problemphyxmini0385-pressureinkilopascals}
  \lean{PhyXMiniProblems.ProblemPhyXMini0385.pressureInKilopascals}
  \uses{def:physics:phyx-mini-0385:phyxminiproblems-problemphyxmini0385-pressureinpascals}
  Kilopascal readout used by the answer choices
\end{definition}
\begin{proof}
  This is the declared model definition of pressure In Kilopascals.
\end{proof}
\begin{definition}[density In Kilograms Per Cubic Meter]
  \label{def:physics:phyx-mini-0385:phyxminiproblems-problemphyxmini0385-densityinkilogramspercubicmeter}
  \lean{PhyXMiniProblems.ProblemPhyXMini0385.densityInKilogramsPerCubicMeter}
  \uses{def:physics:phyx-mini-0385:phyxminiproblems-problemphyxmini0385-massdensityquantity, def:physics:phyx-mini-0385:phyxminiproblems-problemphyxmini0385-densityreadout}
  Kilogram-per-cubic-meter readout of a mass density
\end{definition}
\begin{proof}
  This is the declared model definition of density In Kilograms Per Cubic Meter.
\end{proof}
\begin{definition}[acceleration In Meters Per Second Squared]
  \label{def:physics:phyx-mini-0385:phyxminiproblems-problemphyxmini0385-accelerationinmeterspersecondsquared}
  \lean{PhyXMiniProblems.ProblemPhyXMini0385.accelerationInMetersPerSecondSquared}
  \uses{def:physics:phyx-mini-0385:phyxminiproblems-problemphyxmini0385-accelerationquantity, def:physics:phyx-mini-0385:phyxminiproblems-problemphyxmini0385-accelerationreadout}
  Meter-per-second-squared readout of an acceleration
\end{definition}
\begin{proof}
  This is the declared model definition of acceleration In Meters Per Second Squared.
\end{proof}
\begin{definition}[Cylinder Layer]
  \label{def:physics:phyx-mini-0385:phyxminiproblems-problemphyxmini0385-cylinderlayer}
  \lean{PhyXMiniProblems.ProblemPhyXMini0385.CylinderLayer}
  The three material layers in physical bottom-to-top order
\end{definition}
\begin{proof}
  This is the declared model definition of Cylinder Layer.
\end{proof}
\begin{definition}[Vertical Location]
  \label{def:physics:phyx-mini-0385:phyxminiproblems-problemphyxmini0385-verticallocation}
  \lean{PhyXMiniProblems.ProblemPhyXMini0385.VerticalLocation}
  Distinguished vertical locations needed by the pressure model
\end{definition}
\begin{proof}
  This is the declared model definition of Vertical Location.
\end{proof}
\begin{definition}[Container Material]
  \label{def:physics:phyx-mini-0385:phyxminiproblems-problemphyxmini0385-containermaterial}
  \lean{PhyXMiniProblems.ProblemPhyXMini0385.ContainerMaterial}
  Material named for the container wall
\end{definition}
\begin{proof}
  This is the declared model definition of Container Material.
\end{proof}
\begin{definition}[Cylinder Orientation]
  \label{def:physics:phyx-mini-0385:phyxminiproblems-problemphyxmini0385-cylinderorientation}
  \lean{PhyXMiniProblems.ProblemPhyXMini0385.CylinderOrientation}
  Orientation of the cylinder axis
\end{definition}
\begin{proof}
  This is the declared model definition of Cylinder Orientation.
\end{proof}
\begin{definition}[Figure Label]
  \label{def:physics:phyx-mini-0385:phyxminiproblems-problemphyxmini0385-figurelabel}
  \lean{PhyXMiniProblems.ProblemPhyXMini0385.FigureLabel}
  Literal material and pressure labels visible in the supplied bitmap
\end{definition}
\begin{proof}
  This is the declared model definition of Figure Label.
\end{proof}
\begin{definition}[Figure Height Mark]
  \label{def:physics:phyx-mini-0385:phyxminiproblems-problemphyxmini0385-figureheightmark}
  \lean{PhyXMiniProblems.ProblemPhyXMini0385.FigureHeightMark}
  The height annotations represented by the prose and figure
\end{definition}
\begin{proof}
  This is the declared model definition of Figure Height Mark.
\end{proof}
\begin{definition}[Layered Cylinder Figure]
  \label{def:physics:phyx-mini-0385:phyxminiproblems-problemphyxmini0385-layeredcylinderfigure}
  \lean{PhyXMiniProblems.ProblemPhyXMini0385.LayeredCylinderFigure}
  \uses{def:physics:phyx-mini-0385:phyxminiproblems-problemphyxmini0385-lengthquantity, def:physics:phyx-mini-0385:phyxminiproblems-problemphyxmini0385-cylinderlayer, def:physics:phyx-mini-0385:phyxminiproblems-problemphyxmini0385-verticallocation, def:physics:phyx-mini-0385:phyxminiproblems-problemphyxmini0385-figurelabel, def:physics:phyx-mini-0385:phyxminiproblems-problemphyxmini0385-figureheightmark}
  Independent evidence read from the primary figure. Indices 0, 1, and 2 mean bottom, middle, and top. No pressure answer is stored here
\end{definition}
\begin{proof}
  This is the declared model definition of Layered Cylinder Figure.
\end{proof}
\begin{definition}[Layered Cylinder Setup]
  \label{def:physics:phyx-mini-0385:phyxminiproblems-problemphyxmini0385-layeredcylindersetup}
  \lean{PhyXMiniProblems.ProblemPhyXMini0385.LayeredCylinderSetup}
  \uses{def:physics:phyx-mini-0385:phyxminiproblems-problemphyxmini0385-lengthquantity, def:physics:phyx-mini-0385:phyxminiproblems-problemphyxmini0385-areaquantity, def:physics:phyx-mini-0385:phyxminiproblems-problemphyxmini0385-massdensityquantity, def:physics:phyx-mini-0385:phyxminiproblems-problemphyxmini0385-accelerationquantity, def:physics:phyx-mini-0385:phyxminiproblems-problemphyxmini0385-cylinderlayer, def:physics:phyx-mini-0385:phyxminiproblems-problemphyxmini0385-verticallocation, def:physics:phyx-mini-0385:phyxminiproblems-problemphyxmini0385-containermaterial, def:physics:phyx-mini-0385:phyxminiproblems-problemphyxmini0385-cylinderorientation, def:physics:phyx-mini-0385:phyxminiproblems-problemphyxmini0385-layeredcylinderfigure}
  The physical observables and model parameters. Pressures at distinguished points and throughout the water are stored independently; neither is defined from an answer choice or desired numerical result
\end{definition}
\begin{proof}
  This is the declared model definition of Layered Cylinder Setup.
\end{proof}
\begin{definition}[Matches Layered Cylinder Scenario]
  \label{def:physics:phyx-mini-0385:phyxminiproblems-problemphyxmini0385-matcheslayeredcylinderscenario}
  \lean{PhyXMiniProblems.ProblemPhyXMini0385.MatchesLayeredCylinderScenario}
  \uses{def:physics:phyx-mini-0385:phyxminiproblems-problemphyxmini0385-layeredcylindersetup}
  Qualitative facts stated in the problem prose
\end{definition}
\begin{proof}
  This is the declared model definition of Matches Layered Cylinder Scenario.
\end{proof}
\begin{definition}[Matches Problem Readouts]
  \label{def:physics:phyx-mini-0385:phyxminiproblems-problemphyxmini0385-matchesproblemreadouts}
  \lean{PhyXMiniProblems.ProblemPhyXMini0385.MatchesProblemReadouts}
  \uses{def:physics:phyx-mini-0385:phyxminiproblems-problemphyxmini0385-lengthinmeters, def:physics:phyx-mini-0385:phyxminiproblems-problemphyxmini0385-areainsquaremeters, def:physics:phyx-mini-0385:phyxminiproblems-problemphyxmini0385-pressureinpascals, def:physics:phyx-mini-0385:phyxminiproblems-problemphyxmini0385-layeredcylindersetup}
  Numerical values stated in the prose, excluding the requested pressure
\end{definition}
\begin{proof}
  This is the declared model definition of Matches Problem Readouts.
\end{proof}
\begin{definition}[Matches Supplied Figure]
  \label{def:physics:phyx-mini-0385:phyxminiproblems-problemphyxmini0385-matchessuppliedfigure}
  \lean{PhyXMiniProblems.ProblemPhyXMini0385.MatchesSuppliedFigure}
  \uses{def:physics:phyx-mini-0385:phyxminiproblems-problemphyxmini0385-lengthinmeters, def:physics:phyx-mini-0385:phyxminiproblems-problemphyxmini0385-figurelabel, def:physics:phyx-mini-0385:phyxminiproblems-problemphyxmini0385-layeredcylindersetup}
  Primary-raster evidence: water, gasoline, and air appear bottom to top; the height annotations denote the setup quantities; and P₀ is at the open top. The fill equation also records the unmarked air layer inside the 2.5 m cylinder
\end{definition}
\begin{proof}
  This is the declared model definition of Matches Supplied Figure.
\end{proof}
\begin{definition}[Uses Standard Density And Gravity Calibration]
  \label{def:physics:phyx-mini-0385:phyxminiproblems-problemphyxmini0385-usesstandarddensityandgravitycalibration}
  \lean{PhyXMiniProblems.ProblemPhyXMini0385.UsesStandardDensityAndGravityCalibration}
  \uses{def:physics:phyx-mini-0385:phyxminiproblems-problemphyxmini0385-densityinkilogramspercubicmeter, def:physics:phyx-mini-0385:phyxminiproblems-problemphyxmini0385-accelerationinmeterspersecondsquared, def:physics:phyx-mini-0385:phyxminiproblems-problemphyxmini0385-layeredcylindersetup}
  A separately exposed textbook calibration: water has density 1000 kg/m³, gasoline 750 kg/m³, and gravity is 9.8 m/s². The problem source does not state these three values, so the source-level target below does not take this structure as a premise. It is used only by the calibrated multiple-choice corollary
\end{definition}
\begin{proof}
  This is the declared model definition of Uses Standard Density And Gravity Calibration.
\end{proof}
\begin{definition}[Has Physical Layer Parameters]
  \label{def:physics:phyx-mini-0385:phyxminiproblems-problemphyxmini0385-hasphysicallayerparameters}
  \lean{PhyXMiniProblems.ProblemPhyXMini0385.HasPhysicalLayerParameters}
  \uses{def:physics:phyx-mini-0385:phyxminiproblems-problemphyxmini0385-lengthinmeters, def:physics:phyx-mini-0385:phyxminiproblems-problemphyxmini0385-areainsquaremeters, def:physics:phyx-mini-0385:phyxminiproblems-problemphyxmini0385-pressureinpascals, def:physics:phyx-mini-0385:phyxminiproblems-problemphyxmini0385-densityinkilogramspercubicmeter, def:physics:phyx-mini-0385:phyxminiproblems-problemphyxmini0385-accelerationinmeterspersecondsquared, def:physics:phyx-mini-0385:phyxminiproblems-problemphyxmini0385-cylinderlayer, def:physics:phyx-mini-0385:phyxminiproblems-problemphyxmini0385-layeredcylindersetup}
  Positivity needed for the hydrostatic maximum argument
\end{definition}
\begin{proof}
  This is the declared model definition of Has Physical Layer Parameters.
\end{proof}
\begin{definition}[Satisfies Layered Hydrostatic Laws]
  \label{def:physics:phyx-mini-0385:phyxminiproblems-problemphyxmini0385-satisfieslayeredhydrostaticlaws}
  \lean{PhyXMiniProblems.ProblemPhyXMini0385.SatisfiesLayeredHydrostaticLaws}
  \uses{def:physics:phyx-mini-0385:phyxminiproblems-problemphyxmini0385-lengthquantity, def:physics:phyx-mini-0385:phyxminiproblems-problemphyxmini0385-lengthreadout, def:physics:phyx-mini-0385:phyxminiproblems-problemphyxmini0385-pressurereadout, def:physics:phyx-mini-0385:phyxminiproblems-problemphyxmini0385-densityreadout, def:physics:phyx-mini-0385:phyxminiproblems-problemphyxmini0385-accelerationreadout, def:physics:phyx-mini-0385:phyxminiproblems-problemphyxmini0385-layeredcylindersetup}
  The layered hydrostatic model in arbitrary coherent base units: the open top is at atmospheric pressure; air-weight variation is neglected between the top and gasoline surface; pressure rises by ρ g h through the gasoline; pressure at water depth d rises by ρ\_water g d from the interface; full water depth is the cylinder bottom. No numeric bottom pressure or selected answer occurs in these laws
\end{definition}
\begin{proof}
  This is the declared model definition of Satisfies Layered Hydrostatic Laws.
\end{proof}
\begin{definition}[Is Highest Pressure In Water]
  \label{def:physics:phyx-mini-0385:phyxminiproblems-problemphyxmini0385-ishighestpressureinwater}
  \lean{PhyXMiniProblems.ProblemPhyXMini0385.IsHighestPressureInWater}
  \uses{def:physics:phyx-mini-0385:phyxminiproblems-problemphyxmini0385-lengthquantity, def:physics:phyx-mini-0385:phyxminiproblems-problemphyxmini0385-lengthinmeters, def:physics:phyx-mini-0385:phyxminiproblems-problemphyxmini0385-pressureinpascals, def:physics:phyx-mini-0385:phyxminiproblems-problemphyxmini0385-layeredcylindersetup}
  A candidate is a highest absolute water pressure if it occurs at a physical depth between interface and bottom and dominates every other such depth. The predicate does not select a depth or assign a numerical value by definition
\end{definition}
\begin{proof}
  This is the declared model definition of Is Highest Pressure In Water.
\end{proof}
\begin{definition}[Answer Choice]
  \label{def:physics:phyx-mini-0385:phyxminiproblems-problemphyxmini0385-answerchoice}
  \lean{PhyXMiniProblems.ProblemPhyXMini0385.AnswerChoice}
  Labels of the four displayed answer choices
\end{definition}
\begin{proof}
  This is the declared model definition of Answer Choice.
\end{proof}
\begin{definition}[displayed Pressure In Kilopascals]
  \label{def:physics:phyx-mini-0385:phyxminiproblems-problemphyxmini0385-displayedpressureinkilopascals}
  \lean{PhyXMiniProblems.ProblemPhyXMini0385.displayedPressureInKilopascals}
  \uses{def:physics:phyx-mini-0385:phyxminiproblems-problemphyxmini0385-answerchoice}
  Absolute pressure in kilopascals printed beside each answer label
\end{definition}
\begin{proof}
  This is the declared model definition of displayed Pressure In Kilopascals.
\end{proof}
\begin{definition}[recorded Dataset Answer]
  \label{def:physics:phyx-mini-0385:phyxminiproblems-problemphyxmini0385-recordeddatasetanswer}
  \lean{PhyXMiniProblems.ProblemPhyXMini0385.recordedDatasetAnswer}
  \uses{def:physics:phyx-mini-0385:phyxminiproblems-problemphyxmini0385-answerchoice}
  Dataset answer metadata, deliberately not used as a premise
\end{definition}
\begin{proof}
  This is the declared model definition of recorded Dataset Answer.
\end{proof}
\begin{definition}[displayed Pressure Error]
  \label{def:physics:phyx-mini-0385:phyxminiproblems-problemphyxmini0385-displayedpressureerror}
  \lean{PhyXMiniProblems.ProblemPhyXMini0385.displayedPressureError}
  \uses{def:physics:phyx-mini-0385:phyxminiproblems-problemphyxmini0385-pressureinkilopascals, def:physics:phyx-mini-0385:phyxminiproblems-problemphyxmini0385-answerchoice, def:physics:phyx-mini-0385:phyxminiproblems-problemphyxmini0385-displayedpressureinkilopascals}
  Absolute discrepancy between a physical pressure and a displayed choice, measured in kilopascals
\end{definition}
\begin{proof}
  This is the declared model definition of displayed Pressure Error.
\end{proof}
\begin{definition}[Is Unique Closest Displayed Pressure]
  \label{def:physics:phyx-mini-0385:phyxminiproblems-problemphyxmini0385-isuniqueclosestdisplayedpressure}
  \lean{PhyXMiniProblems.ProblemPhyXMini0385.IsUniqueClosestDisplayedPressure}
  \uses{def:physics:phyx-mini-0385:phyxminiproblems-problemphyxmini0385-answerchoice, def:physics:phyx-mini-0385:phyxminiproblems-problemphyxmini0385-displayedpressureerror}
  A choice is uniquely closest to a physical pressure among the four displayed values. This does not impose a decimal tie-breaking convention: under the chosen calibration the exact 113.25 kPa lies 0.05 kPa from the printed 113.2 kPa, and choice B is nevertheless strictly closer than A, C, or D
\end{definition}
\begin{proof}
  This is the declared model definition of Is Unique Closest Displayed Pressure.
\end{proof}
\begin{lemma}[bottom Pressure Is Highest In Water]
  \label{lem:physics:phyx-mini-0385:phyxminiproblems-problemphyxmini0385-bottompressureishighestinwater}
  \lean{PhyXMiniProblems.ProblemPhyXMini0385.bottomPressureIsHighestInWater}
  \uses{def:physics:phyx-mini-0385:phyxminiproblems-problemphyxmini0385-layeredcylindersetup, def:physics:phyx-mini-0385:phyxminiproblems-problemphyxmini0385-hasphysicallayerparameters, def:physics:phyx-mini-0385:phyxminiproblems-problemphyxmini0385-satisfieslayeredhydrostaticlaws, def:physics:phyx-mini-0385:phyxminiproblems-problemphyxmini0385-ishighestpressureinwater}
  The hydrostatic laws make the bottom pressure maximal throughout water
\end{lemma}
\begin{proof}
  The conclusion follows from the governing relations and figure readouts named in the statement of bottom Pressure Is Highest In Water.
\end{proof}
\begin{lemma}[bottom Pressure In Pascals source Formula]
  \label{lem:physics:phyx-mini-0385:phyxminiproblems-problemphyxmini0385-bottompressureinpascals-sourceformula}
  \lean{PhyXMiniProblems.ProblemPhyXMini0385.bottomPressureInPascals_sourceFormula}
  \uses{def:physics:phyx-mini-0385:phyxminiproblems-problemphyxmini0385-lengthinmeters, def:physics:phyx-mini-0385:phyxminiproblems-problemphyxmini0385-pressureinpascals, def:physics:phyx-mini-0385:phyxminiproblems-problemphyxmini0385-densityinkilogramspercubicmeter, def:physics:phyx-mini-0385:phyxminiproblems-problemphyxmini0385-accelerationinmeterspersecondsquared, def:physics:phyx-mini-0385:phyxminiproblems-problemphyxmini0385-layeredcylindersetup, def:physics:phyx-mini-0385:phyxminiproblems-problemphyxmini0385-satisfieslayeredhydrostaticlaws}
  The strongest numerical relation supported by the source without supplying unstated material data: the bottom pressure is atmospheric pressure plus the gasoline and water hydrostatic heads. Every scalar is a readout of a typed physical quantity in the named coherent SI unit
\end{lemma}
\begin{proof}
  The conclusion follows from the governing relations and figure readouts named in the statement of bottom Pressure In Pascals source Formula.
\end{proof}
\begin{lemma}[bottom Pressure In Kilopascals eq]
  \label{lem:physics:phyx-mini-0385:phyxminiproblems-problemphyxmini0385-bottompressureinkilopascals-eq}
  \lean{PhyXMiniProblems.ProblemPhyXMini0385.bottomPressureInKilopascals_eq}
  \uses{def:physics:phyx-mini-0385:phyxminiproblems-problemphyxmini0385-pressureinkilopascals, def:physics:phyx-mini-0385:phyxminiproblems-problemphyxmini0385-layeredcylindersetup, def:physics:phyx-mini-0385:phyxminiproblems-problemphyxmini0385-matchesproblemreadouts, def:physics:phyx-mini-0385:phyxminiproblems-problemphyxmini0385-usesstandarddensityandgravitycalibration, def:physics:phyx-mini-0385:phyxminiproblems-problemphyxmini0385-satisfieslayeredhydrostaticlaws}
  With the source readouts and the separately stated textbook calibration, the unrounded bottom pressure is 113.25 kPa
\end{lemma}
\begin{proof}
  The conclusion follows from the governing relations and figure readouts named in the statement of bottom Pressure In Kilopascals eq.
\end{proof}
\begin{theorem}[highest Pressure In Water matches choice B under Calibration]
  \label{thm:physics:phyx-mini-0385:phyxminiproblems-problemphyxmini0385-highestpressureinwater-matches-choiceb-undercalibration}
  \lean{PhyXMiniProblems.ProblemPhyXMini0385.highestPressureInWater_matches_choiceB_underCalibration}
  \uses{def:physics:phyx-mini-0385:phyxminiproblems-problemphyxmini0385-pressureinkilopascals, def:physics:phyx-mini-0385:phyxminiproblems-problemphyxmini0385-layeredcylindersetup, def:physics:phyx-mini-0385:phyxminiproblems-problemphyxmini0385-matchesproblemreadouts, def:physics:phyx-mini-0385:phyxminiproblems-problemphyxmini0385-usesstandarddensityandgravitycalibration, def:physics:phyx-mini-0385:phyxminiproblems-problemphyxmini0385-hasphysicallayerparameters, def:physics:phyx-mini-0385:phyxminiproblems-problemphyxmini0385-satisfieslayeredhydrostaticlaws, def:physics:phyx-mini-0385:phyxminiproblems-problemphyxmini0385-ishighestpressureinwater, def:physics:phyx-mini-0385:phyxminiproblems-problemphyxmini0385-isuniqueclosestdisplayedpressure}
  Under the separately stated textbook calibration, the bottom is the highest-pressure point in the water, its unrounded absolute pressure is 113.25 kPa, and B (113.2 kPa) is the unique closest displayed choice
\end{theorem}
\begin{proof}
  The conclusion follows from the governing relations and figure readouts named in the statement of highest Pressure In Water matches choice B under Calibration.
\end{proof}
% --- Archon declaration topology end ---
% --- Archon physics formalization source end ---
